% ============================================================
% 190_T0_Korrekturen_De_ch.tex
% Korrekturverzeichnis – FFGFT / T0-Zeit-Masse-Dualität
% Ältere Dokumente bleiben unverändert; dieses Dokument
% stellt Korrekturen und Präzisierungen klar.
% Stand: 5. Mai 2026 (mit P7, P8: Begrifflichkeit "Renormierung")
% ============================================================

\section*{Vorbemerkung}

Dieses Dokument listet Korrekturen und Präzisierungen zu bereits
veröffentlichten Dokumenten der FFGFT-Serie. Ältere Dokumente werden
\emph{nicht} verändert. Wo ein Widerspruch zwischen einem älteren und
einem neueren Dokument besteht, gilt die hier festgehaltene Version als
verbindlich.

Jeder Eintrag enthält: betroffenes Dokument und Stelle, den fehlerhaften
Ausdruck, den korrekten Ausdruck sowie eine kurze Begründung.

% ============================================================
\section{\texorpdfstring{Korrektur~1: Vorfaktor in der Higgs-Ableitung von $\xi$}{Korrektur 1: Vorfaktor in der Higgs-Ableitung von xi}}
% ============================================================

\subsection*{Betroffenes Dokument}

Dok.~049 (T0-Lagrangian und Higgs-Integration), HTML-Zwischenversion
sowie ältere Arbeitsversionen.

\subsection*{Fehlerhafter Ausdruck}

\begin{equation}
  \xi \;=\; \frac{\lambda_h^2\, v^2}{64\pi^4\, m_h^2}
  \;\approx\; 1{,}0 \times 10^{-5}
  \tag{alt -- falsch}
\end{equation}

\subsection*{Korrekter Ausdruck}

\begin{equation}
  \boxed{
    \xi \;=\; \frac{\lambda_h^2\, v^2}{16\pi^3\, m_h^2}
    \;\approx\; 1{,}33 \times 10^{-4}
  }
  \tag{190.1}
\end{equation}

mit den experimentellen Eingaben:
\begin{align*}
  m_h      &= 125{,}25\,\text{GeV} \quad \text{(Higgs-Masse)}, \\
  v        &= 246{,}22\,\text{GeV} \quad \text{(Vakuum-Erwartungswert)}, \\
  \lambda_h &= \frac{m_h^2}{2v^2} \approx 0{,}129
               \quad \text{(Higgs-Selbstkopplung)}.
\end{align*}

\subsection*{Begründung}

Der Vorfaktor $64\pi^4$ entstammt einem Tippfehler in der frühen
HTML-Visualisierung. Er liefert $\xi \approx 10^{-5}$, was um eine
Größenordnung vom geometrisch abgeleiteten Wert
$\xi = 4/3 \cdot 10^{-4}$ abweicht und die Naturkonstanten nicht
reproduziert. Der korrekte Vorfaktor $16\pi^3$ ergibt
$\xi \approx 1{,}33 \times 10^{-4}$, konsistent mit dem unabhängig
bestimmten Wert aus dem Proton-Elektron-Massenverhältnis (Dok.~009).

% ============================================================
\section{\texorpdfstring{Korrektur~2: Tau-Yukawa-Vorfaktor $r_\tau$}{Korrektur 2: Tau-Yukawa-Vorfaktor r-tau}}
% ============================================================

\subsection*{Betroffenes Dokument}

Dok.~116 (Koide-Formel als Konsequenz der T0-Theorie), Tabelle der
Yukawa-Koeffizienten.

\subsection*{Fehlerhafter Ausdruck}

\begin{equation}
  r_\tau \;=\; \frac{8}{3} \;\approx\; 2{,}667
  \quad \Rightarrow \quad
  m_\tau \approx 1712\,\text{MeV} \quad (-3{,}6\,\% \text{ Abw.})
  \tag{alt -- falsch}
\end{equation}

\subsection*{Korrekter Ausdruck}

\begin{equation}
  \boxed{r_\tau \;=\; \frac{25}{9} \;\approx\; 2{,}778}
  \quad \Rightarrow \quad
  m_\tau \approx 1783\,\text{MeV} \quad (+0{,}4\,\% \text{ Abw.})
  \tag{190.2}
\end{equation}

Die Leptonmassen ergeben sich allgemein aus:
\begin{equation}
  m_\ell \;=\; r_\ell \cdot \xi^{p_\ell} \cdot v
  \tag{190.3}
\end{equation}
mit den Exponenten $p_e = 3/2$, $p_\mu = 1$, $p_\tau = 2/3$
(Schrittweite $\Delta p = 1/3$) und dem Vakuumerwartungswert $v$.

\subsection*{Begründung}

$r_\tau = 8/3$ war ein Überbleibsel aus einer früheren Iterationsstufe.
Dok.~158 (Koide-zu-g-2-Brücke) verwendet bereits $r_\tau = 25/9$, das
numerisch mit der Tau-Masse übereinstimmt. Zudem hat $25/9 = (5/3)^2$
eine geometrische Interpretation als Quadrat der reinen Sexte in der
Naturtonreihe, konsistent mit der musikalischen Resonanz-Analogie
(Dok.~060).

\subsubsection*{Vollständige korrigierte Vorfaktor-Tabelle}

\begin{center}
{\footnotesize
\begin{tabular}{lccc}
\hline
\textbf{Lepton} & \textbf{Exponent $p$} & \textbf{Vorfaktor $r$} & \textbf{$m$ [MeV]}\\
\hline
Elektron   & $3/2$  & $4/3$   & $0{,}511$ \\
Myon       & $1$    & $16/5$  & $105{,}7$ \\
Tau        & $2/3$  & $25/9$  & $1783$    \\
\hline
\end{tabular}
} % footnotesize
\end{center}

% ============================================================
\section{\texorpdfstring{Korrektur~3: Basisformel für $g$-$2$ des Elektrons}{Korrektur 3: Basisformel fuer g-2 des Elektrons}}
% ============================================================

\subsection*{Betroffenes Dokument}

T0-Explorer (HTML-Visualisierung), Erstversion; ältere Zwischenversionen
von Dok.~018.

\subsection*{Fehlerhafter Ausdruck}

\begin{equation}
  a_e \;=\; \frac{4\pi}{f \cdot k_{\text{geom}}}
  \tag{alt -- falsch}
\end{equation}

\subsection*{Korrekter Ausdruck}

\begin{equation}
  \boxed{
    a_e \;=\; \frac{S_3}{f \cdot k_{\text{geom}}}
    \;=\; \frac{2\pi^2}{f \cdot k_{\text{geom}}}
  }
  \tag{190.4}
\end{equation}

mit $f = 1/\xi \approx 7500$, $k_{\text{geom}} = (2/\sqrt{\varphi})\cdot\sqrt{2} \approx 2{,}224$
und $S_3 = 2\pi^2 \approx 19{,}74$ (Oberfläche der 4D-Einheitskugel,
d.\,h.\ des 3D-Randes des Torus).

\subsection*{Begründung}

Der frühere Vorfaktor $4\pi$ entspricht der Oberfläche der 2D-Kugel im
3D-Raum, nicht der geometrisch korrekten 3D-Oberfläche des 4D-Torus.
$S_3 = 2\pi^2$ ist der korrekte Wert aus der Torus-Geometrie. Die
fraktalen Korrekturen für Myon und Tau verwenden weiterhin $4\pi$:

\begin{align}
  \Delta a_\mu  &= \frac{4\pi}{f^{5/3}}, \tag{190.5}\\
  \Delta a_\tau &= \frac{4\pi}{f^{4/3}}. \tag{190.6}
\end{align}

% ============================================================
\section{Präzisierung~1: Genauigkeit der Koide-Formel}
% ============================================================

\subsection*{Betroffene Dokumente}

Dok.~116 (Abstract, Z.\,14): $\Delta Q < 0{,}00003\,\%$ \\
Dok.~158 (Abstract): \enquote{experimentell bestätigt auf $0{,}001\,\%$}

\subsection*{Korrekte Einordnung}

Die Koide-Relation
\begin{equation}
  Q \;=\; \frac{m_e + m_\mu + m_\tau}
               {(\sqrt{m_e}+\sqrt{m_\mu}+\sqrt{m_\tau})^2}
  \;=\; \frac{2}{3}
  \tag{190.7}
\end{equation}
ist experimentell auf besser als $0{,}001\,\%$ bestätigt
(PDG-Werte 2022). Die Angabe $< 0{,}00003\,\%$ in Dok.~116 bezieht
sich auf die interne Konsistenz der T0-Formeln, nicht auf die
experimentelle Messunsicherheit. Beide Aussagen sind inhaltlich korrekt,
aber unterschiedliche Größen. In externen Kommunikationen ist
$0{,}001\,\%$ die geeignete Angabe.

\subsection*{Hinweis}

Die T0-Einzelmassen stimmen auf $\sim 1\,\%$ mit dem Experiment überein;
die hohe Präzision von $Q = 2/3$ entsteht durch die spezifische
algebraische Struktur der Koide-Kombination, nicht durch individuelle
Massengenauigkeit.

% ============================================================
\section{\texorpdfstring{Präzisierung~2: $\hbar$ aus dem Higgs-Feld -- Ableitungskette}{Praezisierung 2: hbar aus dem Higgs-Feld -- Ableitungskette}}
% ============================================================

\subsection*{Kontext}

Die Planck-Konstante $\hbar$ wird in verschiedenen Dokumenten
unterschiedlich hergeleitet. Hier wird die vollständige, konsistente
Ableitungskette festgehalten.

\subsection*{Verbindliche Ableitungskette}

\begin{equation}
  \xi \;=\; \frac{\lambda_h^2 v^2}{16\pi^3 m_h^2}
  \;\longrightarrow\;
  L_0 \;=\; \xi \cdot \ell_P
  \;\longrightarrow\;
  \hbar \;\sim\; \sqrt{\xi}\cdot \ell_P \cdot c
  \tag{190.8}
\end{equation}

Die Zeit-Masse-Bindungsbedingung liefert die physikalische Begründung:
\begin{equation}
  T(x,t)\cdot m(x,t) \;=\; 1
  \quad\Longrightarrow\quad
  T \;=\; \frac{\hbar}{m\, c^2}
  \quad\Longrightarrow\quad
  E \;=\; \hbar\omega.
  \tag{190.9}
\end{equation}

$h = 2\pi\hbar$ ist damit \emph{keine} freie Konstante, sondern eine
geometrische Konsequenz aus $\xi$ und $\ell_P$. Der numerische SI-Wert
$h \approx 6{,}626 \times 10^{-34}\,\text{J·s}$ ist das
Kalibrierungsergebnis dieser Beziehung auf die SI-Basiseinheiten.

% ============================================================
\section{\texorpdfstring{Präzisierung~3: Zwei $\xi$-Parameter -- zwei verschiedene Räume}{Praezisierung 3: Zwei xi-Parameter -- zwei verschiedene Raeume}}
% ============================================================

\subsection*{Kontext}

In den Dokumenten 173--176 (Quantenpunkte, Spin-Qubits, Qubit-Zustandsräume,
Shor-Algorithmus) tauchen zwei numerisch verschiedene $\xi$-Werte auf.
Frühere Dokumente verwendeten diese nicht konsequent getrennt, was zu
Verwirrung führen kann.

\subsection*{Verbindliche Unterscheidung}

\begin{center}
{\footnotesize
\begin{tabular}{p{2.4cm}p{2.0cm}p{2.4cm}p{3.0cm}}
\toprule
\textbf{Parameter} & \textbf{Wert} & \textbf{Geltungs\-bereich} & \textbf{Bedeutung} \\
\midrule
$\xi_0$ \newline (geometrisch) &
$\dfrac{4}{30000}$ \newline $\approx 1{,}33 \times 10^{-4}$ &
3D-geometrischer Raum &
Längen\-hierarchie des Torus; gilt im physikalischen Ortsraum \\[6pt]
$\xi_{\text{Higgs}}$ \newline (algebraisch) &
$\approx 1{,}04 \times 10^{-5}$ &
Zahlen-/Quanten\-zustandsraum &
Algebraische Feldkorrekturen; Higgs-Sektor; Dekohärenzraten \\
\bottomrule
\end{tabular}
}
\end{center}

\subsection*{Begründung}

$\xi_0 = 4/30000$ folgt aus der Torusgeometrie des physikalischen Raums
(geometrische Ableitung, unabhängig von SI-Messungen). $\xi_{\text{Higgs}}$
folgt aus der algebraischen Higgs-Feldstruktur und tritt in Kopplungsformeln
zwischen Quantenzuständen auf (z.\,B.\ Bell-Korrelationen, $T_2$-Hierarchie).

\textbf{Die beiden $\xi$-Werte sind kein Widerspruch}, sondern beschreiben
komplementäre Aspekte: Geometrie des Raums ($\xi_0$) vs.\ Algebra des
Zustandsraums ($\xi_{\text{Higgs}}$). In externen Kommunikationen ist immer
klarzustellen, welcher Parameter gemeint ist.

% ============================================================
\section{\texorpdfstring{Präzisierung~4: $L_0$ und die Schwarzschild-Interpretation}{Praezisierung 4: L0 und die Schwarzschild-Interpretation}}
% ============================================================

\subsection*{Betroffene Dokumente}

Ältere Versionen von Dok.~180--182 (Lagrangian-Ableitung, Torus-Geometrie,
maximale Universum-Skala).

\subsection*{Fehlerhafter Ansatz}

Frühere Versionen interpretierten $L_0 = \xi_0 \cdot \ell_P$ als
Schwarzschild-Radius eines fundamentalen Minimalgebildes. Diese
Interpretation wurde in Dok.~182 (Neufassung 2026) explizit verworfen.

\subsection*{Korrekte Einordnung}

$L_0$ ist ausschließlich die \textbf{geometrische Fundamentallänge} des
FFGFT-Torus:
\begin{equation}
  L_0 \;=\; \xi_0 \cdot \ell_P \;\approx\; 2{,}15 \times 10^{-39}\,\text{m}.
  \tag{190.10}
\end{equation}
Sie definiert die unterste Skalenstufe der $\xi$-Hierarchie und hat
\emph{keine} Schwarzschild-Interpretation. Der kosmologische Hubble-Radius
$R_H$ ist die systemspezifische Obergrenze des kosmologischen Sektors --
nicht eine universelle maximale Größe. Jede Skala hat ihre eigene
systemspezifische Obergrenze.

\subsection*{Anmerkung: Schwarzschild als Konsistenzprüfung}

Die Schwarzschild-Verbindung bleibt als \textbf{interne Konsistenzprüfung} gültig:
Ein Objekt mit Masse $M_0 = \xi_0 \cdot m_P / 2$ hätte den Schwarzschild-Radius
genau $L_0$ -- ohne freien Parameter. Stimmt $M_0$ mit einer Skalenstufe der
$\xi$-Hierarchie überein, ist das ein nichttrivialer Selbstkonsistenztest.
Die Verbindung ist \emph{Probe, nicht Ableitung}.

% ============================================================
\section{Präzisierung~5: Anspruch und Reichweite der Massenformeln}
% ============================================================

\subsection*{Kontext}

Mehrere FFGFT-Dokumente beschreiben Massen-Herleitungen für Leptonen
und andere Teilchen. Der Anspruch dieser Herleitungen muss klar
kommuniziert werden.

\subsection*{Verbindliche Formulierung}

Die FFGFT-Massenformeln (z.\,B.\ $m_\ell = r_\ell \cdot \xi^{p_\ell} \cdot v$)
liefern Werte, die auf $\sim 1\,\%$ mit dem Experiment übereinstimmen.
Dies ist \textbf{keine Präzisionsrechnung}, sondern ein Strukturbeweis:

\begin{itemize}
  \item Die Massenformeln \emph{folgen geometrisch} aus $\xi$ und $v$
    -- sie sind nicht gefittet.
  \item Die verbleibenden $\sim 1\,\%$ Abweichungen spiegeln
    Messunsicherheiten und Näherungen wider, nicht Fehler der Struktur.
  \item Numerisch bessere Berechnungen als das Standardmodell werden
    \emph{nicht} beansprucht.
\end{itemize}

In externen Texten ist zu vermeiden: \enquote{FFGFT berechnet die Massen exakt.}
Korrekt ist: \enquote{FFGFT \emph{leitet} die Massenstruktur aus einem einzigen
Parameter $\xi$ her und reproduziert die Größenordnungen auf $\sim 1\,\%$.}

% ============================================================
\section{Präzisierung~6: Koide-Formel nur mit Originalwerten}
% ============================================================

\subsection*{Verbindliche Regel}

Die Koide-Relation
\begin{equation}
  Q = \frac{m_e + m_\mu + m_\tau}{(\sqrt{m_e}+\sqrt{m_\mu}+\sqrt{m_\tau})^2} = \frac{2}{3}
  \tag{190.11}
\end{equation}
darf ausschließlich mit experimentell gemessenen Massen (PDG-Werte) ausgewertet werden.

\textbf{Nicht zulässig:} Die symbolischen FFGFT-Terme $r_\ell \cdot \xi^{p_\ell} \cdot v$
direkt in die Formel einsetzen. Die $\sqrt{m_i}$-Operation erzeugt Kreuzterme
mit gemischten $\xi$-Potenzen, die dem Torus-Strukturraum fremd sind.
Jede solche Einsetzung ergibt $Q \approx 0{,}6677 \neq 2/3$ und ist
kein Widerspruch zur Theorie, sondern ein Kategorienfehler.

\textbf{Zulässig:} FFGFT-Terme numerisch auswerten
(z.\,B.\ $m_\tau = 1783\,\text{MeV}$) und dieses Ergebnis
wie einen Messwert in die Koide-Formel einsetzen.

\textbf{Mathematischer Hintergrund:} Die $r_\ell$-Terme sind
geometrische Invarianten ihrer Torusmoden. Die Menge der zulässigen
Paare $(r_i, p_i)$ ist bezüglich der Multiplikation nicht abgeschlossen:
$r_e \cdot r_\mu = 64/15$ gehört zu keiner Torusmode.
Das ist das Nichtabschluss-Theorem (Dok.~193).
Die allgemeine Analyse der Kompositionsgrenzen findet sich in Dok.~192.

% ============================================================
\section{\texorpdfstring{Präzisierung~7: Begrifflichkeit \enquote{Renormierungsgruppe} und Skalenstruktur}{Praezisierung 7: Begrifflichkeit Renormierungsgruppe und Skalenstruktur}}
% ============================================================

\subsection*{Betroffene Dokumente}

\begin{itemize}
  \item Dok.~005 (T0-tm-Erweiterung): \enquote{RG-Fluss}, \enquote{Renormierungs\-gruppen-Faktor},
        \enquote{Renormierungs\-gruppen-Invarianz},
        Subsection \enquote{Renormierungs\-gruppen-Behandlung}.
  \item Dok.~008 ($\xi$ und $e$): Subsection \enquote{Modifizierte
        Renormierungs\-gruppen\-gleichungen}.
  \item Dok.~019 (T0-Lagrangian, ältere Fassung):
        Subsection \enquote{Renormierungs\-gruppen-Gleichungen}.
  \item Dok.~086 (Dokumentenübersicht): \enquote{Renormierungsgruppe als Fluss
        in Parameterraum}.
  \item Dok.~189 (Torus-Ableitung): \enquote{kumulativer
        Renormierungs\-gruppen-Lauf}.
  \item Dok.~202 (FFGFT-Feldtheorie Gesamt), \S 18:
        Section-Titel \enquote{Die Renormierungsgruppe}.
\end{itemize}

\subsection*{Bisherige Formulierung}

In den genannten Dokumenten wird der $\xi_0$-modifizierte Skalenlauf
der Kopplung
\begin{equation}
  \frac{d\alpha}{d\ln\mu}
  = \beta_0\,\alpha^2\!\left(1 + \xi_0\,\ln\tfrac{\mu}{E_0}\right)
  \tag{alt -- ungenau}
\end{equation}
und die zugehörige Skalenabhängigkeit als \emph{Renormierungsgruppe}
bezeichnet, mitunter in unmittelbarer Nachbarschaft zu Aussagen, die
$\xi_0$ explizit als topologisch fixierte geometrische Konstante
charakterisieren.

\subsection*{Präzisierte Formulierung}

In FFGFT ist die Skalenabhängigkeit der Kopplung \emph{kein} Resultat
eines Renormierungsverfahrens, sondern eine direkte Konsequenz der
Rekursion $\mathcal{R}$ (Dok.~203) auf der fraktalen Torusgeometrie.
Die korrekte Bezeichnung ist daher:

\begin{equation}
  \boxed{\;
  \begin{aligned}
    &\text{Skalenstruktur aus der Rekursion} \\
    &\quad\equiv\; \text{rekursiver $\alpha$-Lauf via $\mathcal{R}$}
  \end{aligned}
  \;}
  \tag{190.12}
\end{equation}

Konkret:
\begin{itemize}
  \item Der Term $\xi_0\,\ln(\mu/E_0)$ entsteht aus der iterierten
        Anwendung der Rekursion $\mathcal{R}$ auf die Modenstruktur,
        nicht aus einer Counter-Term-Subtraktion.
  \item $\xi_0$ ist \emph{kein} Fixpunkt einer $\beta$-Funktion,
        sondern eine topologisch fixierte Konstante
        ($\xi_0 = 4/30000$, Dok.~180).
  \item Stabilität, Skalenfluss und UV-Endlichkeit sind drei Aspekte
        \emph{derselben} rekursiven Struktur -- keine separaten
        Mechanismen.
\end{itemize}

\subsubsection*{Sprachregelung}

\begin{center}
{\footnotesize
\begin{tabular}{p{4.0cm}p{4.5cm}}
\toprule
\textbf{Vermeiden} & \textbf{Verwenden} \\
\midrule
Renormierungs\-gruppe \newline
{\scriptsize(impliziert QFT-Verfahren)} &
  Skalenstruktur aus der Rekursion \\[6pt]
RG-Fluss / RG-Lauf &
  Rekursiver $\alpha$-Lauf; \newline
  $\mathcal{R}$-induzierte Skalenkorrektur \\[6pt]
Renormierungs\-gruppen-Gleichungen &
  Skalengleichungen aus $\mathcal{R}$ \\[6pt]
Renormierungs\-gruppen-Invarianz &
  Rekursive Skaleninvarianz \\[6pt]
Renormierungs\-gruppen-Fixpunkt &
  Fixpunkt der Rekursion: \newline
  $\mathcal{R}(\Phi_*) = \Phi_*$ \\
\bottomrule
\end{tabular}
}
\end{center}

\subsection*{Begründung}

Die Bezeichnung \enquote{Renormierungsgruppe} importiert ein Werkzeug
der Standard-Quantenfeldtheorie, das auf einem völlig anderen Prinzip
beruht: dort werden UV-Divergenzen durch ein Verfahren absorbiert, und
die Parameter \emph{laufen} unter einer formalen Skalentransformation.
In FFGFT existieren weder das Divergenzproblem (vgl.\ Dok.~159,
\enquote{ohne künstliche Renormierung}) noch ein freier Skalenparameter, der
zu absorbieren wäre. Die Skalenabhängigkeit ist eine \emph{strukturelle
Eigenschaft} der Rekursion, kein nachträglicher Anpassungsschritt.

Die FFGFT-native Sicht ist in Dok.~159 (\enquote{Renormierung als Symptom,
nicht als Lösung}) und Dok.~189 (\enquote{geometrisch bestimmt, nicht durch
Renormierung}) bereits korrekt formuliert; diese Präzisierung dient der
einheitlichen Begrifflichkeit über die ältere Lagrangian- und
$\xi$-Herleitungs-Reihe (Dok.~005, 008, 019) sowie über Dok.~202.

Der zugrundeliegende Mechanismus ist in Dok.~203
(Rekursionsoperator $\mathcal{R}$) und Dok.~204
(Fraktale Randbedingung) ausgearbeitet.

% ============================================================
\section{\texorpdfstring{Präzisierung~8: Begrifflichkeit \enquote{Fraktale Renormierung}}{Praezisierung 8: Begrifflichkeit Fraktale Renormierung}}
% ============================================================

\subsection*{Betroffene Dokumente}

Dok.~005, 008, 012, 014 (mit 19~Vorkommen die häufigste Quelle),
041, 060, 133, 141, 159, 181.

\subsection*{Bisherige Formulierung}

\enquote{Fraktale Renormierung} wird als FFGFT-eigener Terminus für die
geometrischen Korrekturfaktoren bei der Umrechnung zwischen
natürlichen und SI-Einheiten sowie für skalenabhängige Korrekturen
auf dem fraktalen Torus verwendet.

\subsection*{Präzisierte Formulierung}

Da das Wort \enquote{Renormierung} auch in zusammengesetzter Form
das QFT-Konnotat einer Subtraktions- bzw.\ Anpassungsprozedur trägt,
wird der Begriff durch kontextspezifische Alternativen ersetzt:

\begin{center}
{\footnotesize
\begin{tabular}{p{4.0cm}p{4.5cm}}
\toprule
\textbf{Kontext} & \textbf{Verbindliche Bezeichnung} \\
\midrule
Einheitenumrechnung \newline
natürlich $\leftrightarrow$ SI \newline
{\scriptsize(insb.\ Dok.~014)} &
  Geometrische Skalenanpassung \\[8pt]
Skalenabhängige Korrekturen \newline auf dem fraktalen Torus \newline
{\scriptsize(Dok.~133, 159, 181)} &
  Fraktale Korrektur \newline
  {\scriptsize (entspricht Titel von Dok.~133)} \\[8pt]
Allgemeine Skalentransformation \newline durch Rekursion &
  $\mathcal{R}$-induzierte \newline
  Skalenkorrektur \\
\bottomrule
\end{tabular}
}
\end{center}

\subsection*{Begründung}

Der Eigenbegriff \enquote{Fraktale Renormierung} sollte das Verfahren
deutlich von der QFT-Renormierung absetzen, übernimmt jedoch das
zentrale Wort und damit unweigerlich die Konnotation einer
\emph{Korrekturprozedur}. In FFGFT handelt es sich nicht um eine
Prozedur, sondern um die direkte numerische Manifestation der
fraktalen Geometrie. Die vorgeschlagenen Alternativen vermeiden
dieses Missverständnis ohne neuen Begriffsbedarf -- beide Wendungen
sind im Korpus bereits etabliert (Dok.~133 \enquote{Fraktale Korrektur};
\enquote{geometrische Skalenanpassung} im selben semantischen Feld
wie Dok.~014).

% ============================================================
\section*{Aktualisierte Zusammenfassung}
% ============================================================

\begin{center}
{\scriptsize
\resizebox{\textwidth}{!}{%
\begin{tabular}{p{0.5cm}p{2.4cm}p{4.4cm}p{4.4cm}}
\toprule
\textbf{Nr.} & \textbf{Dok.} & \textbf{Fehlerhaft / Unklar} & \textbf{Korrekt / Präzisiert} \\
\midrule
K1 & 049 (HTML) &
  $\xi = \lambda_h^2 v^2/(64\pi^4 m_h^2)$ &
  $\xi = \lambda_h^2 v^2/(16\pi^3 m_h^2)$ \\[3pt]
K2 & 116 (Tab.) &
  $r_\tau = 8/3$ &
  $r_\tau = 25/9$ \\[3pt]
K3 & 018 (Explorer) &
  $a_e = 4\pi/(f \cdot k)$ &
  $a_e = 2\pi^2/(f \cdot k)$ \\[3pt]
P1 & 116 / 158 &
  $\Delta Q < 0{,}00003\,\%$ vs.\ $0{,}001\,\%$ &
  Beide korrekt; $0{,}001\,\%$ für extern \\[3pt]
P2 & mehrere &
  $\hbar$ als Postulat &
  $\hbar$ folgt aus $\xi$ via $L_0 = \xi\cdot\ell_P$ \\[3pt]
P3 & 173--176 &
  Ein einheitlicher $\xi$-Wert &
  Zwei Parameter: $\xi_0$ (geom.) und $\xi_{\text{Higgs}}$ (alg.) \\[3pt]
P4 & 180--182 &
  $L_0$ als Schwarzschild-Radius &
  $L_0$ = geom.\ Fundamentallänge \\[3pt]
P5 & mehrere &
  \enquote{FFGFT berechnet Massen exakt} &
  Massenstruktur aus $\xi$; $\sim 1\,\%$ Übereinstimmung \\[3pt]
P6 & 116 / 158 &
  Koide mit FFGFT-Symbolausdrücken &
  Koide nur mit numerischen Werten (Nichtabschluss, Dok.~193) \\[3pt]
P7 & 005, 008, 019, \newline 086, 189, 202 &
  \enquote{Renormierungsgruppe} als FFGFT-Werkzeug &
  Skalenstruktur aus der Rekursion $\mathcal{R}$ \\[3pt]
P8 & 005, 008, 012, \newline 014, 041, 060, \newline 133, 141, 159, 181 &
  \enquote{Fraktale Renormierung} als Eigenbegriff &
  Geometrische Skalen\-anpassung / fraktale Korrektur \\
\bottomrule
\end{tabular}
}
}
\end{center}
